\section{Conclusion}\label{sec:conclusion}

This study asked whether NBA careers, represented as longitudinal PER sequences, fall into discrete trajectory types, and it did so with a protocol built around reproducibility: held-out model selection, seed and subsample stability, resampling-based consensus with a prespecified rule, equal-terms baselines, tests against structureless references and external validation. A Transformer autoencoder reconstructed careers almost losslessly, but the geometry of its latent space depended on the training seed, so that neither a single UMAP+HDBSCAN run nor a multi-seed consensus produced a reproducible partition beyond a coarse level-by-length structure. Simple summary statistics and a resampled trajectory were far more reproducible and explained more of the trajectory-shape variance. Reproducibility, however, was not evidence of discrete types: representations disagreed with one another, and gap, dip and null-consensus tests indicated a continuum rather than separated groups. We therefore describe careers by three continuous coordinates, level, rise-versus-decline and length, and report the eight-type partition of the summary statistics only as a descriptive segmentation of that continuum, strongly related to honours, career value and draft status but not to position, and we withdraw the earlier claim of nineteen discrete career types. For longitudinal performance data of this kind, the appropriate model is a continuous one, and clustering results should be presented as conventions rather than discoveries unless they survive the tests applied here.
